%% ============================================================
%%  Chapter 3 — Literature Review
%%  NOTE: This chapter will be continuously expanded as the
%%  literature review progresses. Sections are pre-structured
%%  to guide the writing process.
%% ============================================================
\chapter{Literature Review}
\label{ch:literature}

\section{Overview}
\label{sec:lit_overview}

This chapter reviews existing research across four interconnected domains that underpin
this project: (1) eBPF-based network security and observability, (2) DDoS detection
using machine learning, (3) graph-based anomaly detection for network security, and
(4) security in Kubernetes and microservice environments. The review identifies
limitations in existing approaches and establishes the specific research gap that
motivates the proposed framework.

The literature was sourced from IEEE Xplore, ACM Digital Library, USENIX, arXiv, and
Google Scholar. Search queries included: \textit{``eBPF security''}, \textit{``eBPF
intrusion detection''}, \textit{``graph neural network DDoS detection''}, \textit{``Kubernetes
network security''}, \textit{``XDP DDoS mitigation''}, and \textit{``graph anomaly detection
network intrusion''}, filtered to publications from 2018--2025.

% ─────────────────────────────────────────────
\section{eBPF-Based Security Solutions}
\label{sec:lit_ebpf}

\subsection{Foundations of eBPF for Security}

The seminal work by McCanne and Jacobson~\cite{mccanne1993bpf} introduced the original
BPF as a packet filtering engine. The extension to eBPF, documented by Starovoitov and
Borkmann~\cite{borkmann2014ebpf}, transformed this into a general kernel programmability
framework.

Gregg's comprehensive treatment~\cite{gregg2019bpf} established eBPF as a production
observability technology, demonstrating its application across performance tracing,
network monitoring, and security auditing. The work highlights eBPF's key properties
relevant to security: zero-copy telemetry, in-kernel processing, and negligible overhead.

\subsection{XDP-Based DDoS Mitigation}

H\o{}iland-J\o{}rgensen et al.~\cite{hoiland2018xdp} introduced XDP as a high-performance
packet processing framework. Their evaluation demonstrated line-rate packet processing
(up to 24.4 Mpps on a 10 GbE NIC) with sub-microsecond per-packet latency, establishing
XDP as a practical DDoS mitigation substrate.

Miano et al.~\cite{miano2019polycube} presented Polycube, a framework for deploying
eBPF-based network services as composable ``cubes''. Their work demonstrated stateful
firewall and rate-limiting cubes deployable in Kubernetes, achieving throughput
preservation under DDoS conditions with significantly lower overhead than iptables-based
solutions.

\textbf{FlowSentryX}~\cite{flowsentryx} is an open-source XDP/eBPF-based DDoS mitigation
system. It implements per-source-IP rate limiting and SYN cookie validation at the XDP
layer, demonstrating practical deployment on bare-metal servers. However, FlowSentryX
operates at the IP level without Kubernetes-awareness, cannot distinguish between traffic
from different pods sharing an IP (as occurs with Kubernetes Services), and applies static
rate limits without anomaly scoring.

\subsection{Runtime Security and System Call Monitoring}

Falco~\cite{falco2023} is an open-source runtime security tool that uses eBPF (and
historically, a kernel module) to intercept system calls and detect anomalous container
behaviour (e.g., unexpected privilege escalation, shell spawning, file access). While
effective for host-based intrusion detection, Falco operates at the syscall level and does
not model network communication graphs.

Bcc-tools and \texttt{bpftrace}~\cite{bpftrace} provide dynamic tracing of kernel
functions, enabling forensic investigation of security incidents. These tools are used
primarily for offline analysis rather than real-time detection.

\subsection{eBPF-Based Access Control: BPF-LSM}

The BPF Linux Security Module (BPF-LSM) interface, introduced in Linux 5.7~\cite{bpflsm},
allows eBPF programs to be attached to LSM hooks (e.g., \texttt{socket\_connect},
\texttt{file\_open}), enabling fine-grained, programmable access control policies.
Singh et al.~\cite{singh2019krsi} proposed KRSI (Kernel Runtime Security
Instrumentation), demonstrating how BPF-LSM can enforce mandatory access control
policies that adapt to runtime system state — a capability not achievable with static
SELinux or AppArmor policies.

\subsection{Limitations of Existing eBPF Security Work}

\begin{itemize}
  \item Existing eBPF security tools (Falco, FlowSentryX, Cilium) provide excellent
  telemetry and enforcement primitives but do not perform \emph{graph-aware} analysis of
  inter-service communication patterns.
  \item XDP-based mitigation is predominantly reactive and rule-based; it lacks the
  ability to detect subtle, behavioural anomalies in service communication.
  \item No existing eBPF security tool models Kubernetes pod communication as a dynamic
  graph for anomaly detection.
\end{itemize}

% ─────────────────────────────────────────────
\section{Machine Learning for DDoS Detection}
\label{sec:lit_ml_ddos}

\subsection{Classical Machine Learning Approaches}

Early ML-based DDoS detection used flow features (packet rate, byte count, inter-arrival
time) with classifiers such as Random Forest~\cite{rf_ddos}, Support Vector
Machines~\cite{svm_ddos}, and Gradient Boosting~\cite{xgboost_ids}. These approaches
achieve high accuracy on benchmark datasets (CICDDoS2019, CICIDS2017) for volumetric
attacks but degrade significantly for low-and-slow attacks whose per-flow statistics
closely resemble legitimate traffic.

\subsection{Deep Learning Approaches}

LUCID~\cite{lucid2020} (Lightweight Deep Learning for DDoS) uses a convolutional neural
network over time-windowed flow statistics. The authors achieve high F1 scores with low
computational overhead, making LUCID a practical baseline for online detection. However,
LUCID treats each flow independently, discarding inter-flow relational structure.

RNN and LSTM-based approaches~\cite{lstm_ddos} capture temporal patterns in traffic
sequences but still operate on flat, per-flow feature vectors. They do not model
inter-service dependencies.

Kitsune~\cite{mirsky2018kitsune} uses an ensemble of autoencoders for unsupervised
network anomaly detection. Each autoencoder is trained on a feature cluster; anomaly
scores are aggregated via a meta-autoencoder. Kitsune provides unsupervised detection
without labelled attack data but, like other flow-based methods, does not leverage
inter-service relational structure.

\subsection{Federated Learning for DDoS}

Recent work has proposed federated DDoS detection to preserve data privacy across
organisational boundaries~\cite{federated_ddos}. While promising for deployment at
scale, these approaches still rely on flat traffic features and do not address the
intra-cluster visibility problem.

\subsection{Limitations}

\begin{itemize}
  \item All reviewed ML approaches analyse traffic as collections of independent flows,
  missing relational context.
  \item Performance on low-and-slow and application-layer attacks is significantly lower
  than for volumetric attacks.
  \item Models trained on standard datasets (CICDDoS2019) do not generalise to
  Kubernetes microservice traffic, which exhibits different topological and temporal
  characteristics.
\end{itemize}

% ─────────────────────────────────────────────
\section{Graph-Based Anomaly Detection for Network Security}
\label{sec:lit_graph}

\subsection{Graph Neural Networks for IDS}

\textbf{E-GraphSAGE}~\cite{lo2022egraphsage} is the closest existing work to this
project. The authors apply GraphSAGE to network flow graphs where nodes represent
hosts (IP addresses) and edges represent flows. Using features from the CICIDS2018
dataset, they achieve competitive detection performance. However, E-GraphSAGE models a
flat \emph{IP-flow graph} rather than a \emph{Kubernetes service topology graph}.
Crucially, it does not model: (1) pod/service identity and Kubernetes metadata,
(2) dynamic graph evolution as pods are created/destroyed, or (3) intra-cluster
East-West traffic captured via eBPF.

\textbf{Euler}~\cite{euler2022} applies dynamic GNNs to lateral movement detection
in enterprise networks, modelling authentication events as a temporal graph. The
proposed framework shares the temporal graph intuition but targets a different threat
model (host compromise vs. DDoS).

\textbf{Tiresias}~\cite{tiresias2018} uses graph-based methods for network-wide
anomaly prediction, but does not employ GNNs and targets traditional (non-Kubernetes)
networks.

\subsection{Graph Autoencoders for Anomaly Detection}

Ding et al.~\cite{ding2019deep} proposed a deep autoencoder framework for graph
anomaly detection that combines attribute and structural reconstruction errors.
Their framework is directly applicable to communication graph modelling, where nodes
(pods) with anomalous feature vectors \emph{or} anomalous connectivity patterns should
both be flagged.

Dominant-Set Clustering with GAE~\cite{gae_anomaly} has been applied to social network
anomaly detection; adaptations to network security graphs represent an underexplored
direction.

\subsection{Temporal Graph Methods}

Temporal Graph Networks (TGN)~\cite{rossi2020tgn} and TGAT~\cite{xu2020tgat} extend GNNs
to continuous-time dynamic graphs, maintaining per-node memory modules updated by
interaction events. These architectures are directly applicable to streaming eBPF flow
data, where flows arrive as a continuous event stream rather than as periodic snapshots.

\subsection{Limitations of Existing Graph-Based Approaches}

\begin{itemize}
  \item No existing graph-based IDS incorporates Kubernetes-native telemetry (pod identity,
  namespace, service topology, Kubernetes metadata) as graph structure.
  \item Existing work uses static datasets (CICIDS) not representative of dynamic
  Kubernetes workloads.
  \item Integration of graph anomaly detection with real-time enforcement (rate limiting,
  NetworkPolicy) has not been demonstrated.
\end{itemize}

% ─────────────────────────────────────────────
\section{Kubernetes and Container Security}
\label{sec:lit_k8s_security}

\subsection{Network Policy Enforcement}

Kubernetes NetworkPolicy~\cite{k8s_networkpolicy} provides a declarative API for
specifying ingress/egress rules. However, enforcement depends on the CNI plugin, and the
API supports only static, predefined policies. Dynamic policy adaptation based on runtime
anomaly scores is not natively supported.

CiliumNetworkPolicy~\cite{cilium_network_policy} extends the standard NetworkPolicy with
Layer 7 (HTTP, DNS) rules and identity-based (rather than IP-based) policies, making it
more suitable for dynamic Kubernetes environments.

\subsection{Microservice Communication Modelling for Fault Detection}

MicroRank~\cite{yu2021microrank} and Microscope~\cite{lin2018microscope} model
Kubernetes microservice communication as call graphs for root-cause analysis of
performance anomalies. These works demonstrate that graph modelling of microservice
interactions is feasible and informative but focus on reliability (latency, errors) rather
than security (attack detection).

\subsection{Container Runtime Security}

Falco~\cite{falco2023} and Tracee~\cite{tracee2022} (Aqua Security) use eBPF/kernel
modules to monitor container syscall behaviour and detect security policy violations.
These solutions complement network-level detection but do not address the DDoS threat
model.

% ─────────────────────────────────────────────
\section{Comparison with Proposed Framework}
\label{sec:lit_comparison}

\begin{table}[h]
\centering
\caption{Comparison of existing solutions with the proposed \sys{} framework}
\label{tab:comparison}
\begin{tabularx}{\textwidth}{lcccccc}
\toprule
\textbf{System} & \textbf{eBPF} & \textbf{K8s-aware} & \textbf{Graph ML} & \textbf{Temporal} & \textbf{Adaptive Mitig.} & \textbf{Low-\&-Slow} \\
\midrule
FlowSentryX    & \checkmark & \texttimes & \texttimes & \texttimes & \texttimes & \texttimes \\
Falco          & \checkmark & \checkmark & \texttimes & \texttimes & \texttimes & \texttimes \\
Kitsune        & \texttimes & \texttimes & \texttimes & \texttimes & \texttimes & Partial \\
LUCID          & \texttimes & \texttimes & \texttimes & \texttimes & \texttimes & \texttimes \\
E-GraphSAGE    & \texttimes & \texttimes & \checkmark & \texttimes & \texttimes & Partial \\
Euler          & \texttimes & \texttimes & \checkmark & \checkmark & \texttimes & Partial \\
\textbf{\sys{}} & \checkmark & \checkmark & \checkmark & \checkmark & \checkmark & \checkmark \\
\bottomrule
\end{tabularx}
\end{table}

\noindent Table~\ref{tab:comparison} illustrates that no existing system combines all five
capabilities simultaneously. This combination --- eBPF telemetry, Kubernetes-native graph
modelling, temporal graph ML, and adaptive mitigation --- constitutes the primary novelty
claim of this project.

% ─────────────────────────────────────────────
\section{Summary and Research Gap}
\label{sec:lit_gap}

The literature review reveals the following converging limitations:

\begin{enumerate}
  \item \textbf{Flat traffic analysis:} Existing ML-based DDoS detection systems treat
  network traffic as a collection of independent flows, discarding the relational structure
  of microservice communication.

  \item \textbf{Perimeter focus:} Industry-grade solutions (Cloudflare, AWS Shield) operate
  at the network edge and have no visibility into intra-cluster, pod-to-pod communication.

  \item \textbf{Static mitigation:} Rate limiting and blocking rules are threshold-based and
  cannot adapt to dynamic, evolving attack patterns without human intervention.

  \item \textbf{Graph ML without eBPF integration:} Graph-based IDS research
  (E-GraphSAGE, Euler) demonstrates the power of relational learning but relies on
  pre-captured, offline datasets rather than real-time eBPF telemetry.

  \item \textbf{eBPF security without graph modelling:} eBPF-based tools (Cilium, Falco,
  FlowSentryX) provide excellent telemetry and enforcement but do not perform
  graph-aware anomaly detection.
\end{enumerate}

\textbf{The gap:} No existing system integrates real-time eBPF telemetry within a
Kubernetes cluster, models pod communication as a dynamic graph, applies temporal GNN-based
anomaly detection, and drives adaptive mitigation from graph anomaly scores. This project
proposes to fill this gap with \sys{}.

\bigskip
\noindent \textit{Note: This literature review will continue to be expanded as additional
papers are reviewed during the project. Estimated final coverage: 40--60 papers across all
reviewed domains.}
